\chapter{Numerical Examples}

\section{Experimental Setting}

This chapter demonstrates lists, subfigures, and numerical-result tables. Replace the sample text and placeholders with actual experimental settings and results.

\begin{itemize}
  \item Dataset or test problem: describe the source and basic scale.
  \item Baseline methods: list the compared methods.
  \item Metrics: describe error, accuracy, time, or other evaluation metrics.
\end{itemize}

\section{Subfigures}

Figure \ref{fig:subfig-example} shows two placeholder subfigures.
\begin{figure}[htbp]
  \centering
  \begin{subfigure}{0.45\linewidth}
    \centering
    \fbox{\rule{0pt}{30mm}\rule{55mm}{0pt}}
    \caption{First result}
    \label{fig:subfig-a}
  \end{subfigure}
  \hfill
  \begin{subfigure}{0.45\linewidth}
    \centering
    \fbox{\rule{0pt}{30mm}\rule{55mm}{0pt}}
    \caption{Second result}
    \label{fig:subfig-b}
  \end{subfigure}
  \caption{A sample subfigure layout}
  \label{fig:subfig-example}
\end{figure}

\section{Result Table}

Table \ref{tab:results} shows a sample result table.
\begin{table}[htbp]
  \centering
  \caption{Sample numerical results}
  \label{tab:results}
  \begin{tabular}{lcc}
    \toprule
    Method & Metric 1 & Metric 2 \\
    \midrule
    Method A & 0.123 & 1.23 \\
    Method B & 0.098 & 1.05 \\
    \bottomrule
  \end{tabular}
\end{table}
